\begin{answers}{ch:lm}{言語モデル：次のトークンを予測する}
\item 損失は 0.462、パープレキシティは 1.587 です。
組は「たけ」「けや」「やぶ」「ぶや」「やけ」「けた」の 6 個です。
回数の表は、た→け 1、け→や 1、け→た 1、や→ぶ 1、や→け 1、ぶ→や 1 です。
確率の表は、た→け と ぶ→や が 1、け と や の行はそれぞれ 2 通りに分かれて 0.5 ずつです。
6 回の予測のうち 4 回が確率 0.5、2 回が確率 1 なので、損失は $\frac{1}{6}(4\log 2+2\cdot 0)=\frac{2}{3}\log 2=0.462$ です。
パープレキシティは $e^{0.462}=2^{2/3}=1.587$ です。
\file{ch06/toy_bigram.py} の \code{text} の行を \code{text = "たけやぶやけた"} にすると、同じ値が出ます。

\item $p_b=C_{ab}/C_a$ とおくと、$\sum_b p_b=1$ なので、$p$ も確率分布です。
$C_{ab}=C_a p_b$ を代入すると、次のようになります。
\[
-\sum_b C_{ab}\log q_b=C_a\Bigl(-\sum_b p_b\log q_b\Bigr)=C_a\Bigl(-\sum_b p_b\log p_b+\sum_b p_b\log\frac{p_b}{q_b}\Bigr)
\]
右辺の第 2 項は $\KL{p}{q}$ です。ギブスの不等式（式\eqref{eq:math-gibbs}）から、これは 0 以上です。
0 になるのは $q=p$ のときだけです。
第 1 項は $q$ によらないので、全体が最小になるのは $q_b=p_b=C_{ab}/C_a$ のときです。
なお、$C_a=0$ の行は損失に入らないので、どんな $q$ でもかまいません。

\item （コード）文章が途中から、ほとんど見たことのない漢字の羅列に変わります。
種 0 では、元の「こう積みんながまま\dots」が「こう積み賜応じまま\dots」となり、後半は「彼女滋肌羃条偕刺戟渓牆督隗デ肛\dots」になりました。
理由は、平滑化で、訓練データにない組にも確率が付くことです。
行 $a$ の「見たことのない次の文字」の確率の合計は、$\alpha(V-\text{現れた種類数})/(C_a+\alpha V)$ です。
「。」の行では 0.003 ですが、762 回しか出ない「吾」の行では 0.35 になります。
訓練データの位置で平均すると 0.12 で、$\alpha=0.01$ の 0.03 の約 4 倍です。
一度まれな文字を引くと、その行はほぼ一様分布です。そのため、次もでたらめな文字を引きやすく、ずれが続きます。
これは top-$k$ の節で見た「すそ」の問題そのものです。1 個 1 個は小さな確率でも、4078 個ぶん足すと無視できません。
損失（採点）には平滑化が必要ですが、生成では top-$k$ や top-$p$ ですそを切るほうがよい、ということです。

\item （コード）学習率 0.05 では 100 回後の訓練損失が 3.688 で、まだ下がりきっていません。
学習率 5 では訓練損失は 3.489 まで下がりますが、検証損失は 4.152 と、0.5 のとき（3.625）よりずっと悪くなりました。
この二つは、\figref{fig:lm-lr-sweep}の左端と右端のグラフです。学習率を 0.2 と 2 にした場合も、同じ図で見比べられます。
Adam は、勾配の大きさによらず、各成分をほぼ学習率くらいずつ動かします。
学習率 5 では、訓練データにない組のロジットが一気に大きく下げられ、検証データの未観測の組に付く確率が小さくなりすぎます。
\file{ch06/sweep_lr.py} の出力では、検証データのうち訓練データにない組の損失が、学習率 0.5 の 15.93 から、5 では 43.22 に増えていました（訓練データにある組は、どちらも約 3.38）。
SGD に替えると、学習率 0.5 では 100 回後も 8.253 で、ほとんど進みません。
勾配は式\eqref{eq:lm-grad}のとおり $C_a/N$ 倍されていて、ほとんどの成分が非常に小さいからです。
学習率を 50 にしても 6.127、500 にしても 4.751 でした。
SGD では、まれな文字の行がなかなか動かないためです。成分ごとに歩幅をそろえる Adam が、この問題では有利です。

\item 1 文字あたりのパープレキシティは $129.4^{1/1.75}=16.1$ です。
パープレキシティは損失の指数なので、割るべきなのは損失です。
$4.863/1.75=2.78$ を指数に戻すと $e^{2.78}=e^{4.863/1.75}=(e^{4.863})^{1/1.75}$ になります。
「1.75 で割る」は、パープレキシティでは「1.75 乗根を取る」に当たります。
単位をそろえずに 129.4 と 29.9 を比べると、「基盤モデルはトライグラムより 4 倍も迷っている」と誤解します。
実際は、1 回の予測で 1.75 文字ぶんを当てているので、問題が難しいだけです。
1 文字あたりにそろえると 16.1 と 29.9 で、基盤モデルのほうがよいのです。

\item （コード）温度 0.7、top-$p$ 0.9 で、種 0, 1, 2 から次の 3 本ができました。
「吾輩は主人は、それがまたことに、誰にもある。それを断って、どうする」、
「吾輩は人の子供は赤い灯影に半ば自然に勝手に、その日の間からである。」、
「吾輩はちょっとも限らないでに観察して見せた。／　その代りにした。／」。
すそが切られるので、意味の通る句が続きやすく、貪欲法のような同じ句の繰り返しもありません。
top-$k$ の $k=1$ は、貪欲法（温度 0）と同じです。
残る候補が最大の確率の 1 個だけなので、くじを引いても必ずそれが選ばれます。
実際、種を変えても、貪欲法と同じ「この時の彼は、その時の彼は、\dots」が出ました。
\end{answers}
